\section{Experimentos de dietas en mosca soldado negra}

Para evaluar el impacto de diferentes dietas en la eficiencia del proceso de bioconversión de residuos orgánicos mediante la mosca soldado negra (\textit{Hermetia illucens}), se realizarán experimentos controlados utilizando desechos orgánicos provenientes de diferentes fuentes, como sustrato principal. Estos residuos incluirán restos de alimentos ricos en materia orgánica, cuya composición permitirá analizar el efecto de distintas mezclas en el crecimiento larval, la tasa de conversión de residuos y la generación de biomasa. Durante el proceso, se implementará un sistema de monitoreo en tiempo real mediante sensores IoT para medir variables ambientales claves como temperatura, humedad, iluminación, pH y concentración de nutrientes esenciales (NPK).
 
Estos parámetros serán fundamentales para evaluar cómo las condiciones ambientales influyen en la eficiencia de la bioconversión y en las emisiones de gases de efecto invernadero (GEI). Los datos recopilados serán analizados  utilizando técnicas estadísticas y algunos de ellos con algoritmos de inteligencia artificial (IA), con el objetivo de maximizar la producción de biomasa y subproductos de valor agregado, al tiempo que se minimizan las emisiones de GEI. Esta estrategia permitirá validar la viabilidad de la mosca soldado negra como solución sostenible para la gestión de residuos alimentarios y contribuirá al diseño de un sistema automatizado de gestión de residuos en entornos agrícolas y agroindustriales.